\chapter{Tag 12 — Reed-Solomon-Prüfbytes für Datamatrix}

\begin{quote}
Heute kommt der Aha-Moment, für den wir die ganze Reed-Solomon-Mathe
aus Etappe 6, 7 und 8 gebaut haben. Wir nehmen unsere fünf
HONIG-Codewort-Bytes von gestern und berechnen die sieben
Prüfbytes, die ein Datamatrix-Decoder zur Fehlerkorrektur braucht.
Der Encoder von Tag 7 funktioniert dafür fast 1:1 — wir müssen nur
\emph{ein} Detail anpassen: das Modulo-Polynom für GF($2^8$).
\end{quote}

\section*{Lernziele für heute}

Am Ende von Tag 12 kannst du \dots

\begin{itemize}
\item erklären, warum Datamatrix ein \emph{anderes} Modulo-Polynom
  als unser Tag-7-Code verwendet — und dass das nur Konvention ist
\item eine neue GF($2^8$)-Klasse aus der alten kopieren und genau
  eine Konstante austauschen
\item das Generator-Polynom für 7 Prüfsymbole im Datamatrix-Stil
  bauen
\item für deine HONIG-Codewörter \texttt{[73, 80, 79, 74, 72]} die
  sieben Prüfbytes berechnen lassen — mit deinem eigenen Code, nicht
  mit einer fertigen Bibliothek
\end{itemize}

\paragraph{Material} Bleistift, Karopapier, Laptop mit Thonny und
deinem Code aus Tag 7 als Referenz.

% =============================================================
\section{Erinnerung: was Reed-Solomon an Tag 7 gemacht hat \normalfont(\textit{≈ 15 min})}
% =============================================================

In Tag 7 haben wir einen Reed-Solomon-Encoder gebaut, der
\textbf{drei Schritte} macht:

\begin{enumerate}
\item \textbf{Generator-Polynom $g(x)$ aufstellen}: ein Produkt von
  Linearfaktoren über GF($2^8$),
  $g(x) = (x - \alpha^1)(x - \alpha^2) \cdots (x - \alpha^t)$,
  mit $\alpha = 2$ als primitivem Element und $t$ als Anzahl der
  Prüfsymbole.
\item \textbf{Daten mit Nullen auffüllen}: an die Datenfolge werden
  $t$ Nullen angehängt — das wird das Eingabe-Polynom $D(x)$.
\item \textbf{Polynomdivision modulo $g(x)$}: wie schriftliche
  Division, nur dass Subtraktion durch XOR ersetzt wird (das ist
  Addition in GF($2^8$)). Der Rest ist die Prüfbyte-Folge.
\end{enumerate}

Dieser Algorithmus ist \emph{universell} — er funktioniert für jeden
endlichen Körper und jede Anzahl Prüfsymbole. Das einzige, was sich
zwischen Anwendungen unterscheidet, sind:

\begin{itemize}
\item das \textbf{Modulo-Polynom}, mit dem GF($2^8$) konstruiert wird
\item die \textbf{Anzahl der Prüfsymbole} $t$ (bestimmt durch die
  Symbol-Größe)
\end{itemize}

\section{Datamatrix nutzt ein anderes Modulo-Polynom \normalfont(\textit{≈ 15 min})}

Du erinnerst dich an Tag 5: GF($2^n$) wird konstruiert, indem man
modulo eines \emph{irreduziblen} Polynoms vom Grad $n$ rechnet. Es
gibt mehrere irreduzible Polynome — welches man wählt, ist Konvention.

Tag 7 hat das \emph{klassische} Reed-Solomon-Polynom verwendet:

\[
  \texttt{0x11D} = x^8 + x^4 + x^3 + x^2 + 1
  \quad (\text{Bitmuster } \texttt{1\,0001\,1101})
\]

\textbf{Datamatrix} (im Standard ISO/IEC 16022) wählt aber ein
anderes:

\[
  \texttt{0x12D} = x^8 + x^5 + x^3 + x^2 + 1
  \quad (\text{Bitmuster } \texttt{1\,0010\,1101})
\]

Beide sind irreduzibel. Beide ergeben Körper mit 256 Elementen.
Beide funktionieren mit $\alpha = 2$ als primitivem Element. Aber:
die \emph{Multiplikationstabellen sind unterschiedlich}, und damit
auch die α-Potenzen, das Generator-Polynom und am Ende die
Prüfbytes. Wenn du den falschen Polynom-Wert nimmst, scannt das
Handy dein Symbol nicht.

\subsection*{Eine ganz neue GF-Klasse?}

Wir müssen nicht alles neu bauen. Wir kopieren die GF256-Klasse aus
Tag 7 und ändern \emph{eine einzige Zeile}: die \texttt{MOD\_POLY}-
Konstante.

\begin{pythoneinheit}{1}{GF256DM — die Datamatrix-Variante}
\pythoncodeteil{tag12/gf256_dm.py}{klasse}

Vergleich zu Tag 7: identisch bis auf die erste Konstante in der
Klasse (\texttt{MOD\_POLY = 0x12D} statt \texttt{0x11D}). Die
Multiplikation, die Potenzierung, die Addition (XOR) — alles
genauso wie zuvor.
\end{pythoneinheit}

\begin{pythoneinheit}{2}{α-Potenzen-Tabelle als Sanity-Check}
\pythoncodeteil{tag12/gf256_dm.py}{alpha}

Wenn $\alpha = 2$ wirklich primitiv ist, müssen $\alpha^0, \alpha^1,
\ldots, \alpha^{254}$ alle 255 Nicht-Null-Elemente von GF($2^8$)
genau einmal liefern, und $\alpha^{255}$ muss wieder $1$ sein.
Genau das prüft das Snippet — und gibt \texttt{True} aus, wenn das
Datamatrix-Polynom in Ordnung ist.
\end{pythoneinheit}

\section{Generator-Polynom für 7 Prüfsymbole \normalfont(\textit{≈ 15 min})}

Für ein 12$\times$12-Datamatrix brauchen wir $t = 7$ Prüfbytes. Das
Generator-Polynom ist also:

\[
  g(x) = (x - \alpha^1)(x - \alpha^2)(x - \alpha^3)(x - \alpha^4)
       (x - \alpha^5)(x - \alpha^6)(x - \alpha^7).
\]

Sieben Linearfaktoren ergeben ein Polynom vom Grad 7 mit 8
Koeffizienten. Wir müssen sie nicht von Hand ausrechnen — die
Funktion aus Tag 7 macht das (jetzt mit der Datamatrix-GF):

\begin{pythoneinheit}{3}{generator\_polynom — übernommen aus Tag 7}
\pythoncodeteil{tag12/rs_dm.py}{generator}

Identisch zur Tag-7-Variante, nur mit \texttt{from gf256\_dm import
\dots}~statt der alten GF-Klasse. Eingabe: Anzahl der Prüfsymbole.
Ausgabe: Liste der Polynom-Koeffizienten als GF256DM-Objekte, vom
höchsten Grad zum niedrigsten.
\end{pythoneinheit}

% =============================================================
\section{Polynomdivision für HONIG \normalfont(\textit{≈ 25 min})}
% =============================================================

\begin{bleistiftuebung}{1}{Polynomdivision von Hand — die ersten zwei Schritte}
Aus Etappe 11 hast du die HONIG-Datenbytes \texttt{[73, 80, 79, 74,
72]}. Wir hängen sieben Nullen an (für die sieben Prüfbytes):

\[
  D = [73,\; 80,\; 79,\; 74,\; 72,\; 0,\; 0,\; 0,\; 0,\; 0,\; 0,\; 0]
\]

Schreibe \emph{nur} den ersten Schritt der Polynomdivision modulo
$g(x)$ auf:

\begin{itemize}
\item Der Leitkoeffizient ist 73 ($\alpha^?$ — schlag in deiner
  GF-Tabelle nach, welche Potenz das ist).
\item Multipliziere $g(x)$ mit diesem Leitkoeffizienten — gibt 8
  GF($2^8$)-Produkte.
\item Subtrahiere (= XOR) das Ergebnis von $D$ — die ersten 8
  Koeffizienten ändern sich, davon der erste wird $0$.
\end{itemize}

Der zweite Schritt geht genauso, mit dem neuen Leitkoeffizienten an
zweiter Stelle. Insgesamt brauchst du \textbf{fünf Schritte}, dann
sind alle Daten-Stellen abgearbeitet, und in den letzten sieben
Stellen steht der Rest — die Prüfbytes.

(Die ganze Polynomdivision von Hand wäre mühsam — fünf Schritte
mit je acht GF($2^8$)-Multiplikationen. Greta soll die ersten zwei
Schritte machen, um das Konzept zu spüren; den Rest übernimmt Python
gleich.)
\end{bleistiftuebung}

\begin{pythoneinheit}{4}{Reed-Solomon-Encoder für Datamatrix}
\pythoncodeteil{tag12/rs_dm.py}{encoder}

Identisch zur Tag-7-Implementierung. Der Algorithmus ist universell;
nur die GF-Klasse ist eine andere.
\end{pythoneinheit}

\begin{pythoneinheit}{5}{Demo: HONIG ECC berechnen}
\pythoncodeteil{tag12/rs_dm.py}{demo}

Ausgabe (kannst du in Thonny direkt ausprobieren):

\begin{minted}{text}
Daten: [73, 80, 79, 74, 72]
ECC : [70, 186, 97, 167, 44, 40, 243]
Komplettes Codewort: [73, 80, 79, 74, 72, 70, 186, 97, 167, 44, 40, 243]
\end{minted}

Zwölf Codewort-Bytes — fünf Daten, sieben Prüfbytes. Genau das, was
in das 12$\times$12-Datamatrix-Symbol hineingehört.
\end{pythoneinheit}

\begin{bleistiftuebung}{2}{GRETA-ECC ausrechnen lassen}
Wenn du gestern die Codewort-Bytes für \texttt{GRETA}
(\texttt{[72, 83, 70, 85, 66]}) ausgerechnet hast, lass deinen
Code jetzt die sieben Prüfbytes dafür ausspucken — du brauchst sie
in Etappe 14, wenn du dein eigenes Symbol zeichnest.
\end{bleistiftuebung}

% =============================================================
\section{Reflexion und Brücke zu Tag 13 \normalfont(\textit{≈ 10 min})}
% =============================================================

Du hast jetzt:

\begin{itemize}
\item eine neue GF($2^8$)-Variante, die Datamatrix versteht
\item das Generator-Polynom mit sieben Wurzeln
\item zwölf Codewort-Bytes für HONIG (5 Daten + 7 Prüf)
\end{itemize}

Was noch fehlt:

\begin{itemize}
\item Der \textbf{Datenplatzierungs-Algorithmus} — wie kommen die
  $12 \cdot 8 = 96$ Bits in die 96 Datenmodule des
  12$\times$12-Symbols. Datamatrix nutzt ein „Treppenmuster" mit
  L-förmigen Codewort-Blöcken, die durch das Symbol diagonal wandern.
  Das ist \textbf{Tag 13}.
\item L-Pattern und Zeitsignal als Frame, der Werkstattbogen, das
  Hand-Zeichnen und der finale Scan-Test — das ist \textbf{Tag 14}.
\end{itemize}

\paragraph{Was du heute eigentlich erlebt hast}
Die Reed-Solomon-Mathematik aus Etappe 6 bis 8 wirkte abstrakt. Heute
hat sie sich \emph{konkret an Datenbytes für ein echtes Symbol}
bewährt. Der Polynom-Wechsel ist ein winziges Detail — der
Algorithmus dahinter bleibt derselbe, von Tag 6 bis hierher. Genau
das macht die Stärke endlicher Körper aus: einmal verstanden, immer
wieder anwendbar, mit minimaler Anpassung pro Standard.
